\section{引言}

\subsection{本讲义的主题} 这是 2015 年 PCMI 研究生暑期学校上一门短期课程的讲义。目标有二。其一，向听众介绍一些有表示论应用的几何对象；Springer 纤维及其推广正是这样的好例子。其二，介绍轨道积分，着重它与仿射 Springer 纤维的联系，从而为 B-C.~Ng\^o 关于基本引理的短期课程提供背景材料。

本讲义的几何部分（除 \S\ref{s:orb} 以外）研究三类「纤维」：Springer 纤维、仿射 Springer 纤维与 Hitchin 纤维，复杂度依次递增。我们会讨论它们的几何性质，例如连通性与不可约分支。我们会在这些簇上构造若干群作用，并借助这些作用研究几个不平凡的例子。最重要的是，我们会研究这些纤维的上同调上的某些 Weyl 群作用，这些作用并不来自簇本身的自同构。三类纤维在表示论上的意义，以及它们之间的类比，可以概括为下表。

\begin{table}[htbp]
\centering
\begin{tabular}{|c|c|c|c|}
\hline
 & Springer 纤维 & 仿射 Springer 纤维 & Hitchin 纤维\\
\hline
域 & $k$ & 局部域 $F=k\lr{t}$ & 整体域 $k(X)$\\
\hline
对称性 & $W$ & $\tilW$ & $\tilW$\\
\hline
「扩张」对称性 & 分次 AHA & 分次 DAHA & 分次 DAHA\\
\hline
表示论 & $G(k)$ 的特征标 & $G(F)$ 的轨道积分 & $k(X)$ 上 $G$ 的迹公式\\
\hline
\end{tabular}
\end{table}

这里 AHA 指仿射 Hecke 代数，DAHA 指双仿射 Hecke 代数；$X$ 是 $k$ 上的一条代数曲线；$W$ 与 $\tilW$ 分别是 Weyl 群与扩张仿射 Weyl 群。

本讲义不追求对所有论断给出完整证明，而是着重指出有趣的现象与例子。不过，对几个关键结果我们给出大体完整的证明，例如
\begin{itemize}
\item 定理 \ref{th:Spr corr}（Springer 对应）；
\item 定理 \ref{th:ft quot}（仿射 Springer 纤维的有限性）；
\item 定理 \ref{th:orb coho}（稳定轨道积分的上同调解释）。
\end{itemize}

\subsection{对读者的预备要求} 本讲义面向对几何表示论感兴趣的研究生新生。我们假定读者了解一些基本的代数几何（概形理论、模问题、有限域上的点计数等），尽管偶尔会用到代数叠的语言。我们也假定读者了解一些 Lie 理论（代数闭域上与局部域上的约化群），但只要熟悉 $\GL_{n}$ 与 $\SL_{n}$，就足以读懂本讲义的大部分内容。

下面说明本讲义使用的上同调理论。我们处理的是任意域 $k$ 上的代数簇，而非 $\CC$ 上的，因此使用系数在 $\ell$ 进层（通常是常层 $\Ql$）中的平展上同调。我们把 $k$ 上概形 $X$ 的 $\Ql$ 系数平展上同调简记为 $\cohog{*}{X}$。不熟悉平展上同调的读者，不妨先取 $k=\CC$，把 $\cohog{*}{X}$ 理解为 $X(\CC)$ 以 $\QQ$ 为系数的奇异上同调。反常层只在 \S\ref{ss:Spr corr} 中用到。

\subsection*{致谢} 感谢 2015 年 PCMI 暑期项目的共同组织者、各位授课人与工作人员。也感谢听讲者提出的反馈。特别感谢陈竞仁（Jingren Chi）仔细通读了本讲义初稿，并提出有益的建议。
